\section{Factoring with verifiable auxiliary data}\label{sec:process}

We next organize the order trial into a factoring process.  Its successful
output contains enough auxiliary information to compute, retrospectively,
the exact probabilities of all its useful transitions.  This information
is produced by the process itself; it is not supplied as advice.
Throughout this section, $n$ and $B=2^n$ retain their values for the original
input $N$, even when we work on a smaller integer.  We use $K$ repetitions
within each randomized transition, at most $L$ transition slots, and $W$
bits per transition output, where
\begin{equation}\label{eq:process-parameters}
 K=n^5,\qquad L=n^{10},\qquad W=n(K+1).
\end{equation}

For $M\ge2$, define $\D(M)$ to be the following rooted occurrence tree.
The root records $M$ and its distinct prime factors in increasing order,
together with their positive exponents.  For each distinct odd prime $p$ in that factorization, it has one
child containing $\D(p-1)$; the children are ordered by $p$.  Equal labels
under different parents are separate occurrences.  There is no child for
$p=2$.  Every descent decreases the label, so this defines unique finite
data.  We write $\D=\D(N)$.
For example, the root of $\D(21)$ has children $\D(2)$ and $\D(6)$,
associated with its prime factors $3$ and $7$. The second child has its
own $\D(2)$ child, associated with the prime factor $3$ of $6$; the two
occurrences of $\D(2)$ remain separate.

The recursive $p-1$ structure recalls Pratt's succinct primality
certificates~\cite{Pratt}. Here there are no primitive-root witness fields:
prime labels are checked by Lemma~\ref{lem:primality}, and the auxiliary
factorizations instead determine exact orders and transition probabilities
retrospectively. The order-to-gcd splitting mechanism follows the
Miller--Shor reduction~\cite{Miller,Shor}; its required unconditional
counting argument is included in the proof below.

\begin{proposition}[A process with verifiable transition probabilities]
\label{prop:process}
There is a uniformly polynomial-time process, using the order trial of
Proposition~\ref{prop:order-trial}, with the following properties.
\begin{enumerate}
\item It either aborts or generates $\D$.  A deterministic polynomial-time
verifier accepts exactly these complete data.  The process's classical controller
makes at most $L$ transitions, and completed runs are padded to exactly $L$
by constant-output dummy transitions.  Each transition output has a fixed
$W$-bit encoding.
\item Given $\D$ and the recorded classical input of a transition, one can
compute a good-output predicate, one canonical good encoding, and the exact
probability $s$ that an ordinary execution of that transition has a good
output.  The probability $s$ is dyadic, has polynomial bit length, and obeys
\[
 s\ge1-2^{-2n}.
\]
For a dummy transition, its constant output is good and canonical, and
$s=1$.
\item Every evolution whose transition outputs are all good completes the
tasks and generates $\D$ without aborting or reaching a resource cap.  This
assertion also holds when good outputs are supplied directly.  The work is
polynomially bounded on arbitrary output sequences as well, with malformed
outputs causing an abort.
\end{enumerate}
Only the transition output is included in its $W$-bit encoding.  Internal
order-trial registers and recorded computational history are separate
workspace and may be retained.
\end{proposition}

The $W$-bit output is the string consumed by the controller. The completed
tree supplies the arithmetic information for later tests, while the
separate history and quantum work registers are retained for the inverse
circuit. These are three different records; the output width does not
bound the whole workspace.

\begin{proof}
We construct the process and verify these properties in turn.
Each transition uses fresh work registers and fresh randomness.  Its
classical output is copied to a history register before the controller
uses it; earlier quantum work is retained as inactive workspace.
Consequently, conditional on any recorded classical prefix, the next
transition has exactly the distribution specified by its recorded
classical inputs, even if earlier work registers remain entangled.

\paragraph{The information contained in $\D$.}
Suppose that a subfactor $m$ is being split while factoring a node $M$.
From $\D$ we can factor $m$ by division from the prime list for $M$.
Moreover, $\D$ supplies the factorization of $p-1$ for every odd prime
$p\mid m$.  Thus it supplies the factorization of
\begin{equation}\label{eq:totient-from-data}
 \varphi(m)=\prod_{p^e\parallel m}(p-1)p^{e-1}.
\end{equation}
Here the formula follows by counting units modulo a prime power and using
the Chinese remainder bijection; the latter follows from B\'ezout's
identity.  The contribution from $p=2$ needs no child.

Given this factorization and a unit $a$ modulo $m$, its order is computable
exactly.  Start with exponent $v=\varphi(m)$.  For each prime $q$ in its known
factorization, repeatedly replace $v$ by $v/q$ while $q\mid v$ and
$a^{v/q}\equiv1\pmod m$.  The order divides
$\varphi(m)$ by the coset proof of Lagrange's theorem.  Throughout the
procedure it divides $v$, and at termination no excess prime factor can
remain.  Hence the final $v$ is the order, and the remaining prime
exponents give its factorization.  There are polynomially many modular
powerings.  The same procedure applies to each prime-power component of
$m$.  These are retrospective computations: they are used to evaluate
recorded transitions once the data have been generated.

We shall use the elementary bound
\begin{equation}\label{eq:totient-lower-bound}
 \frac{\varphi(d)}d\ge\frac1{n+1}\qquad(1\le d<B).
\end{equation}
Indeed, if $p_1<\cdots<p_k$ are the distinct prime divisors of $d$, then
$k\le n$ and $p_i\ge i+1$, so the product formula gives
$\varphi(d)/d\ge\prod_{i=1}^k i/(i+1)=1/(k+1)$.
The empty product covers $d=1$.

\paragraph{The controller and its two transitions.}
To factor a node, apply the deterministic primality test in
Lemma~\ref{lem:primality} to each current factor.  A prime is recorded;
an even composite can be split by $2$; and a composite perfect power can
be split by finding an integral root.  Root searches for exponents at
most $n$ take polynomial time.  Consequently the remaining case is an odd
$m$ having at least two distinct prime divisors.  We use the following two
kinds of transition for this case.

A list output consists of $K$ fields of $b_m=\lceil\log_2m\rceil$ bits,
followed by zero padding to width $W$; an order output consists of one
$n$-bit integer followed by zero padding to width $W$.  A dummy output is
the all-zero $W$-bit string.  The controller rejects malformed padding,
and every good-output predicate below requires a valid encoding.

First, a \emph{list transition} generates $K$ independent uniform
$b_m$-bit integers.  Entries outside
$[0,m)$ and nonunits are ignored.  A list is good if some usable entry $a$
has unequal values of
\[
 \nu_2\bigl(\ord_{p^e}(a)\bigr),\qquad p^e\parallel m,
\]
where $\nu_2$ denotes the exponent of $2$ in a positive integer.
We now compute its good probability without enumerating residues.

The nonzero elements of a prime field form a cyclic group.  For
completeness, let $T$ be the least common multiple of their orders.
Lagrange's theorem gives $T\mid p-1$, while the root bound for $X^T-1$
gives $p-1\le T$.  The exponent of a finite abelian group is attained:
for each prime power dividing that exponent, take a suitable power of an
element whose order contains that prime power, then multiply these
elements of pairwise coprime orders.  Thus there is an element of order
$T=p-1$.

Write $p-1=h2^\ell$, with $h$ odd.  In this cyclic group the counts at
levels $0,1,\ldots,\ell$ are
$h,h,2h,\ldots,2^{\ell-1}h$.  Reduction of units modulo $p^e$ to units
modulo $p$ has an odd kernel of size $p^{e-1}$.  The order of an element
and the order of its reduction therefore have the same $2$-adic level,
and each residue has $p^{e-1}$ lifts.  The corresponding counts modulo
$p^e$ are
\begin{equation}\label{eq:level-counts}
 p^{e-1}h,\quad p^{e-1}h,\quad
 2p^{e-1}h,\quad\ldots,\quad 2^{\ell-1}p^{e-1}h.
\end{equation}
No level has more than half the units.

Let $c_i(t)$ be these counts for the prime-power components of $m$,
extended by zero outside their ranges.  The Chinese remainder bijection
makes the levels independent for a uniform unit.  Since there are at
least two components, the probability that all levels agree is at most
$1/2$: for the first two level laws $q_1,q_2$,
$\sum_t q_1(t)q_2(t)\le\max_t q_1(t)\le1/2$.
The number of favorable residues is therefore exactly
\[
 F=\varphi(m)-\sum_t\prod_i c_i(t)\ge\frac{\varphi(m)}2.
\]
There are at most $n$ components and $n+1$ levels, so this formula is
polynomial-time computable from $\D$.  By
\eqref{eq:totient-lower-bound} and $2^{b_m}\le2m$, the list's exact good
probability is
\begin{equation}\label{eq:list-good-probability}
 s=1-\left(1-\frac{F}{2^{b_m}}\right)^K,
 \qquad \frac{F}{2^{b_m}}\ge\frac1{4(n+1)}.
\end{equation}
A canonical good list has first entry equal to the Chinese remainder
residue that is $1$ on the component corresponding to the smallest prime
and $-1$ on all others, followed by zeros.  Its first entry has levels $0$ and $1$.

Second, for every usable entry $a$, an \emph{order transition} runs $K$
independent order trials.  It outputs the smallest returned candidate,
or $0$ if there was none.  If $r=\ord_m(a)$, goodness means that this
output is $r$, which is also the canonical good output.  By
Proposition~\ref{prop:order-trial}, every nonempty candidate is a positive
multiple of $r$, and the probability of the candidate $r$ in one trial is
$\varphi(r)\lambda(r)$, where
$\lambda(r)=2^{-\lceil\log_2r\rceil-10}$.  Consequently
\begin{equation}\label{eq:order-good-probability}
 s=1-\bigl(1-\varphi(r)\lambda(r)\bigr)^K,
 \qquad \varphi(r)\lambda(r)\ge\frac1{2048(n+1)}.
\end{equation}
The last bound follows from \eqref{eq:totient-lower-bound} and
$2^{\lceil\log_2r\rceil}\le2r$, also when $r=1$.

Equations~\eqref{eq:list-good-probability} and
\eqref{eq:order-good-probability} give exactly computable dyadics.
Their denominator exponents are at most $nK$ and $(n+10)K$, respectively.
For either transition, the lower bound on the one-trial success parameter
$x$ gives
\[
 Kx\ge\frac{n^5}{2048(n+1)}
       \ge\frac{n^4}{4096}\ge2n
       \qquad(n\ge128).
\]
Thus $(1-x)^K\le e^{-2n}<2^{-2n}$, proving the required uniform bound on
$s$.  The good predicates, probabilities, and canonical outputs are all
computable from $\D$ and the classical inputs $m$ or $(m,a)$.

\paragraph{Why good transitions produce a split.}
For every positive even order output $r'$, try
$\gcd(a^{r'/2}-1,m)$ and accept it only if it is a proper nontrivial
divisor.  If no entry produces a split, abort.
If the list and its order transitions are good, choose an entry with
unequal local levels and let $r$ be its true order.  By the Chinese
remainder bijection,
\[
 r=\operatorname*{lcm}_{p^e\parallel m}\ord_{p^e}(a).
\]
The largest local $2$-adic level is positive.  The local orders at smaller
levels divide $r/2$, whereas
at a largest-level component $a^{r/2}$ has order two.  The only square
roots of $1$ modulo an odd prime power are $1$ and $-1$: in
$(x-1)(x+1)\equiv0\pmod{p^e}$ the entire odd prime power divides one
factor.  Hence $a^{r/2}$ equals $1$ on the smaller-level components and
$-1$ on the largest-level components.  The gcd is a proper nontrivial
divisor.  This proves that all-good outputs prevent an abort at every
split, regardless of how those outputs were obtained.

After completely factoring a node, the controller creates its specified
children and factors them in a fixed order.  Once all these tasks have
finished, it fills unused transition slots with dummy transitions.

\paragraph{Size bounds and verification.}
The binary splitting tree for one node $M$ has at most $\log_2M$ leaves:
all leaf labels are at least two and their product is $M$.  A partial run
has at most $\log_2M$ split attempts as well, including a possible final
failed attempt.  At each auxiliary-tree descent,
\[
 \sum_{\substack{p\mid M\\p>2}}\log_2(p-1)
 \le\sum_{p\mid M}\log_2p\le\log_2M.
\]
Thus the total logarithmic size at any depth is at most $\log_2N<n$.
In two successive descents $M\to p-1\to q-1$, the intermediate label is
even and $q>2$ divides it, so $q-1<(p-1)/2<M/2$.
There are therefore fewer than $2n+2$ levels.  Since every node label is
at least two, each contributes at least one to the logarithmic size;
the tree has at most $(2n+2)n$ node occurrences.  Summing the split-attempt
bound over the same levels gives at most $(2n+2)n$ split tasks.
Each uses at most $K+1$ transitions, whence
\[
 \#\{\text{list and order transitions}\}
 \le(2n+2)n(K+1)\le6n^7<n^{10}=L.
\]
All deterministic work and all $K$ trials within an order transition
also have polynomial bounds.

These bounds hold even for bad output sequences.  Only genuine gcd
splits are accepted, prime leaves are verified, and children are created
only from the resulting actual prime factors.  An abort only shortens
the work.  Fixed polynomial caps on work and storage can therefore be
imposed without excluding any all-good evolution.  Each node's prime
list has at most $n$ entries of polynomial bit length, so the occurrence
bound also bounds the storage for the complete data.  A list has at most
$nK$ unpadded bits, and an order value, including zero, has $n$ bits.
Thus the specified encodings fit within $W$ bits and give each canonical
output a unique string.
The controller depends only on these classical outputs, and internal
quantum work can be retained without being treated as part of an output.

Finally, the complete data are polynomially verifiable.  Require the root
label to be $N$. Before arithmetic, reject node labels outside $[2,B)$,
proposed prime labels outside $[2,M]$ at a node $M$, exponents outside
$[1,n]$, and lists or trees exceeding the preceding polynomial size caps.
At a node $M$, require strictly increasing prime labels, check primality
using Lemma~\ref{lem:primality}, and check the factorization's product.
Then require exactly the ordered children $p-1$ for the distinct odd primes
in that list and apply the same
checks recursively.  Reject missing or extra children and malformed or
oversized records; a partial product may be rejected as soon as it exceeds
$M$. These guards make the verification polynomial even on arbitrary
invalid records, and no true data violate them.
Unique prime factorization implies that an accepted tree is exactly
$\D(N)$; only the operational split history may vary.  The all-good
splitting argument ensures that every all-good evolution supplies this
complete tree.  This proves the proposition.
\end{proof}
